Courier operations that coordinate parcels manually or with disconnected tools give customers little visibility, leave responsibility for each parcel unclear and keep no reliable record of status changes. This thesis designs, implements and evaluates the Smart Courier and Parcel Management System (SCPMS), a web application that gives one courier operator a shared, role-controlled record of every parcel and pushes changes to authorised users in real time.

SCPMS is built on MongoDB, Express, React and Node.js, with Socket.IO messaging and OpenStreetMap-based maps. Customers register with e-mail verification, book parcels and follow them on a map. Delivery agents confirm pickups by scanning the customer's QR code, advance parcels through a server-enforced six-state lifecycle and share their location. Administrators assign agents, manage accounts and export CSV and PDF reports, and anyone with a tracking code sees a reduced public view. Security measures include short-lived access tokens with rotating, reuse-detecting refresh cookies, role and ownership checks, account lockout, rate limiting and NoSQL-injection filtering.

A new black-box harness ran 66 API and Socket.IO test cases against the unmodified backend and an in-memory database; all 66 passed. Production dependencies have no known vulnerabilities. The evaluation also records fourteen findings, including 98 frontend lint errors, a permissive forward-skip rule in the lifecycle and a QR step that identifies parcels without proving handover. No usability study, load test or field trial was performed. The work shows that a secure, real-time courier system can be built with open-source technologies; usability and operational benefit remain to be established.

\vspace{1em}
\noindent\textbf{Keywords:} courier management system; parcel tracking; real-time web applications; Socket.IO; role-based access control; MERN stack; web application security; software testing
